\section{The level-two Rademacher expansion}\label{sec:rademacher}

This section states the exact formula for the coefficients of the finite part on which the proof of Theorem~\ref{thm:all} rests: a vector-valued Rademacher series over the three cosets of $\Gamma^0(2)$, whose principal parts at the two cusps are the wall sums of Section~\ref{sec:walls} and whose correction terms come from the shadow. The next four sections prove it.

\subsection{Conventions, frames and Weil sums}\label{circ:conv}
For $w\in\CC\setminus(-\infty,0]$ and $s\in\RR$ we put $w^s=\exp(s\Log w)$ with the principal logarithm, $-\pi<\Arg w<\pi$. All powers below are of numbers in $\HH$, in the lower half-plane, or in the right half-plane, so they are always defined, and $w^sw^{s'}=w^{s+s'}$.

Let $\Gamma^0(2)$ be the group of matrices of $SL_2(\ZZ)$ with even upper right entry. It has index three, and
\[
R_I=\begin{pmatrix}1&0\\0&1\end{pmatrix},\qquad R_T=\begin{pmatrix}1&1\\0&1\end{pmatrix},\qquad R_S=\begin{pmatrix}0&-1\\1&0\end{pmatrix}
\]
represent its right cosets. For $M\in SL_2(\ZZ)$ the \emph{frame} of $M$ is the $F\in\{I,T,S\}$ with $MR_F^{-1}\in\Gamma^0(2)$. The frames $I$ and $T$, which differ by a half period, belong to the cusp $\infty$, and the frame $S$ to the cusp $0$.\footnote{The duality group $\Gamma_1(2)$ acts on the charges; the finite part $\varphi^F_m(\sigma,v)$ is a function of $\sigma$ on which the relevant group, after the substitution $x=\e(\sigma/2)$, is $\Gamma^0(2)$. The cusp $\infty$ has width two, which is why it carries two frames.}

Fix $m\ge1$. For $\gamma\ge1$ let $\mathcal C_\gamma$ be the set of pairs $(\gamma,\delta)$ with $0\le-\delta\le2\gamma-1$ and $\gcd(\gamma,\delta)=1$. For such a pair let $\alpha\in\{0,\dots,\gamma-1\}$ be the inverse of $\delta$ modulo $\gamma$ ($\alpha=0$ if $\gamma=1$), $\beta=(\alpha\delta-1)/\gamma$, and
\[
M_{\gamma,\delta}=\begin{pmatrix}\alpha&\beta\\\gamma&\delta\end{pmatrix},\qquad N_{\gamma,\delta}=M_{\gamma,\delta}^{-1}=\begin{pmatrix}\delta&-\beta\\-\gamma&\alpha\end{pmatrix},\qquad N^*_{\gamma,\delta}=\begin{pmatrix}\delta&\beta\\\gamma&\alpha\end{pmatrix},
\]
all in $SL_2(\ZZ)$, with $j_{\gamma,\delta}(\tau)=-\gamma\tau+\alpha$ the automorphy factor of $N_{\gamma,\delta}$, and let $F(\gamma,\delta)$ be the frame of $N_{\gamma,\delta}$. For $\tilde\ell,\ell\in\ZZ$ put
\begin{equation}\label{eq:weil}
W_{\gamma,\delta}(\tilde\ell,\ell)=\frac1{\sqrt{2m\gamma i}}\sum_{T=0}^{\gamma-1}\e\Big(\frac{\alpha x^2-2\ell x+\delta\ell^2}{4m\gamma}\Big),\qquad x=\tilde\ell-2mT .
\end{equation}
This is the matrix coefficient of the Weil representation attached to the lattice $\sqrt{2m}\,\ZZ$, evaluated at $N_{\gamma,\delta}$ \cite[(2.41)]{CNR21}; it depends on $\tilde\ell$ only modulo $2m$ (Lemma~\ref{circ:lem:weil}).

For $\ell\in\ZZ$ let $\vartheta_\ell(\tau)=\sum_{r\equiv\ell\,(2m)}\e(r^2\tau/4m)$ be the unary theta series of weight $\frac12$ and index $m$, so that $\vartheta_\ell(\tau)=\vartheta_{m,\ell}(\tau,0)$ for the theta functions $\vartheta_{m,\ell}(\tau,z)=\sum_{r\equiv\ell\,(2m)}\e(r^2\tau/4m+rz)$ of \cite[\S5]{EZ85}. Its coefficients are real, so
\begin{equation}\label{circ:thetaconj}
\vartheta_\ell(-\bar w)=\overline{\vartheta_\ell(w)}\qquad(w\in\HH).
\end{equation}
With $N=N_{\gamma,\delta}$, $N^*=N^*_{\gamma,\delta}$ and $j=j_{\gamma,\delta}$, for $w,\tau\in\HH$,
\begin{equation}\label{circ:Nstar}
N(-\bar w)=-\overline{N^*w},\qquad N^*w+N\tau=\frac{w+\tau}{(\gamma w+\alpha)\,j(\tau)} .
\end{equation}
The first identity is immediate; for the second, the numerator $(\delta w+\beta)(-\gamma\tau+\alpha)+(\delta\tau-\beta)(\gamma w+\alpha)$ collapses to $(\alpha\delta-\beta\gamma)(w+\tau)=w+\tau$.

\begin{lemma}\label{lem:weilbound}
$|W_{1,\delta}|=(2m)^{-1/2}$, and $|W_{\gamma,\delta}(\tilde\ell,\ell)|\le\sqrt{\gcd(m,\gamma)/m}$ for every $\gamma\ge1$.
\end{lemma}

\begin{proof}
For $\gamma=1$ the sum has one term. In general, $\alpha x^2/(4m\gamma)=\alpha\tilde\ell^2/(4m\gamma)-\alpha\tilde\ell T/\gamma+\alpha mT^2/\gamma$, so up to a factor of modulus one the sum is $\sum_{T\bmod\gamma}\e\big((\alpha mT^2+bT)/\gamma\big)$ for an integer $b$. A quadratic Gauss sum modulo $\gamma$ with leading coefficient $\alpha m$, $\gcd(\alpha,\gamma)=1$, has modulus at most $\sqrt{2\gcd(m,\gamma)\gamma}$ \cite[Chapter~1]{BEW98}. Dividing by $\sqrt{2m\gamma}$ gives the claim.
\end{proof}

\subsection{Polar data and shadows}
For $m\ge1$ the polar terms in each frame are the pairs $(\tilde\nu,\tilde\ell)$ with $0\le\tilde\ell<2m$ and $\tilde D=4m\tilde\nu-\tilde\ell^2<0$, where $\tilde\nu\in\ZZ$ in the frame $S$ and $\tilde\nu\in\frac12\ZZ$ in the frames $I$ and $T$. Their coefficients are the wall sums of Section~\ref{sec:walls}:
\begin{equation}\label{eq:polardata}
a_S(\tilde\nu,\tilde\ell)\ \text{from \eqref{eq:P0}},\qquad a_I(\tilde\nu,\tilde\ell)\ \text{from \eqref{eq:Pinf}},\qquad a_T(\tilde\nu,\tilde\ell)=(-1)^{2\tilde\nu}a_I(\tilde\nu,\tilde\ell).
\end{equation}
By Lemma~\ref{lem:walls}(iv) a wall with $m_G,n_G\ge-1$ has $t\le t_0=m^2+4m$, so only finitely many polar terms have a nonzero coefficient. That these numbers are the principal parts of $\varphi^F_m$ is the content of two theorems proved below:
\begin{enumerate}
\item[(P$\infty$)] the principal part of $\varphi^F_m$ at the cusp $\infty$ is given by \eqref{eq:Pinf} (Theorem~\ref{thm:Pinf});
\item[(P0)] the principal part of the finite part at the cusp $0$, in the frame $S$, is given by \eqref{eq:P0} (Theorem~\ref{thm:P0}).
\end{enumerate}

The shadow of the finite part is governed by the coefficient $\hat f(m)=d_1(m)$ of the polar part \eqref{eq:polar}. In the three frames it enters with the constants and orders
\begin{equation}\label{eq:mockdata}
(c_S,\mu_S)=\big(\tfrac{1}{16}\hat f(m),\,1\big),\qquad (c_I,\mu_I)=\big(\hat f(m),\,\tfrac12\big),\qquad (c_T,\mu_T)=\big(-\hat f(m),\,\tfrac12\big).
\end{equation}
The number $\mu_F$ is the order of the pole of $1/g(\sigma/2)$ in the local parameter of the frame: $1/g$ has the term $x^{-1}=\e(-\sigma/2)$ at $\infty$, and its transform under $\sigma\mapsto-1/\sigma$ starts with $\e(-\sigma)/16$, by $\eta(-1/\tau)=\sqrt{-i\tau}\,\eta(\tau)$.\footnote{With $\tau=\sigma/2$, the substitution $\sigma\mapsto-1/\sigma$ is $\tau\mapsto-1/4\tau$, and $g(-1/4\tau)=\eta(-1/4\tau)^8\eta(-1/2\tau)^8=2^{12}\tau^8\eta(\sigma)^8\eta(2\sigma)^8$, whose expansion starts with $\e(\sigma)$.} These are the leading coefficients of the functions $G_F$ of Proposition~\ref{tr:prop:jacobi} (Remark~\ref{tr:rem:G}). For $\gamma\ge1$ and $g\in\ZZ$ put
\[
f_{\gamma,g}(u)=\begin{cases}\dfrac{\pi^2}{\sinh^2\!\big(\frac{\pi u}{\gamma}-\frac{i\pi g}{2m\gamma}\big)},& 2m\gamma\nmid g,\\[2.2ex]
\dfrac{\pi^2}{\sinh^2(\pi u/\gamma)}-\dfrac{\gamma^2}{u^2},& 2m\gamma\mid g,\end{cases}\qquad u\in\RR,
\]
the second extended by $f_{\gamma,g}(0)=-\pi^2/3$. Since $\sinh(a-ib)=0$ for real $a,b$ only when $a=0$ and $b\in\pi\ZZ$, $f_{\gamma,g}$ is continuous on $\RR$; it depends on $g$ only modulo $2m\gamma$.

\subsection{The expansion}
For a charge $(m,n,\ell)$ with $m\ge1$, $0\le\ell\le m$ and $D=2mn-\ell^2>0$ define
\begin{align}
\mathcal M(m,n,\ell)&=\sum_{\gamma\ge1}\sum_{(\gamma,\delta)\in\mathcal C_\gamma}\ \sum_{(\tilde\nu,\tilde\ell)}a_{F}(\tilde\nu,\tilde\ell)\,\e\Big(\frac{\alpha\tilde D+\delta D}{4m\gamma}\Big)W_{\gamma,\delta}(\tilde\ell,\ell)\,\frac{2\pi}\gamma\Big(\frac{|\tilde D|}{D}\Big)^{\frac{15}4}I_{\frac{15}2}\Big(\frac{\pi\sqrt{|\tilde D|D}}{m\gamma}\Big),\label{eq:radmain}\\
\mathcal K_2(m,n,\ell)&=-\sqrt{2m}\sum_{\gamma\ge1}\gamma^{-\frac12}\!\!\sum_{(\gamma,\delta)\in\mathcal C_\gamma}\!\! c_F\,\e\Big(\frac{-4m\mu_F\alpha+\delta D}{4m\gamma}\Big)W_{\gamma,\delta}(0,\ell)\Big(\frac{4m\mu_F}{D}\Big)^4I_8\Big(\frac{2\pi}\gamma\sqrt{\frac{\mu_FD}m}\Big),\label{eq:radK2}\\
\mathcal K_3(m,n,\ell)&=\frac1{2\pi}\sum_{\gamma\ge1}\gamma^{-2}\!\!\sum_{(\gamma,\delta)\in\mathcal C_\gamma}\ \sum_{g=-m\gamma}^{m\gamma-1}c_F\,\e\Big(\frac{-\alpha(4m\mu_F+g^2)+\delta D}{4m\gamma}\Big)W_{\gamma,\delta}(g,\ell)\Big(\frac{4m\mu_F}{D}\Big)^{\frac{17}4}\label{eq:radK3}\\
&\qquad\times\int_{-\sqrt{\mu_F/m}}^{\sqrt{\mu_F/m}}f_{\gamma,g}(u)\Big(1-\frac{mu^2}{\mu_F}\Big)^{\frac{17}4}I_{\frac{17}2}\Big(\frac{2\pi}\gamma\sqrt{\frac{(\mu_F-mu^2)D}m}\Big)\,du,\notag
\end{align}
where in each line $F=F(\gamma,\delta)$ and $\alpha$ is attached to $(\gamma,\delta)$ as above.

\begin{theorem}[Rademacher expansion]\label{thm:rad}
For every charge $(m,n,\ell)$ with $m\ge1$, $0\le\ell\le m$ and $D=2mn-\ell^2>0$, the three series \eqref{eq:radmain}--\eqref{eq:radK3} converge absolutely, and
\begin{equation}\label{eq:hypR}
(-1)^{\ell+1}c^F_m(n,\ell)=\tfrac12(-1)^\ell\sqrt i\,\big(\mathcal M+\mathcal K_2+\mathcal K_3\big)(m,n,\ell).
\end{equation}
\end{theorem}

This is Theorem~\ref{thm:radB} of the introduction. The proof occupies Sections~\ref{sec:polarinf} to~\ref{circ:sec}. Section~\ref{sec:polarinf} proves (P$\infty$) and Section~\ref{sec:polar0} proves (P0), by following the Fourier coefficients of $1/\Phit_6$ and of a second product $1/\Phi^S_6$ across the Farey walls. Section~\ref{tr:sec} proves the theta decomposition of the finite parts in the three frames, the completion of the Appell--Lerch sum by a Mordell integral and the resulting transformation law. Section~\ref{circ:sec} runs the circle method on Farey arcs and ends with the proof of Theorem~\ref{thm:rad}.\footnote{For weakly holomorphic forms at level $N$ the expansion is the circle method of Rademacher and Zuckerman \cite{RZ38,Zuc39}. For mixed mock Jacobi forms on $SL_2(\ZZ)$, the case of $1/\Phi_{10}$, it was proved by Bringmann and Manschot \cite{BMa13}, with the correction terms coming from the completion of Zwegers \cite{Zwe02}, and it was used for black holes by Ferrari and Reys \cite{FR17}. At level two the expansion is written down, and assumed, in the derivation of the immortal index by Cardoso, Nampuri and Rossell\'o \cite[\S4]{CNR21}.}

\begin{lemma}\label{lem:radconv}
The three series \eqref{eq:radmain}--\eqref{eq:radK3} converge absolutely.
\end{lemma}

\begin{proof}
There are at most $2\gamma$ pairs in $\mathcal C_\gamma$, the polar terms are finitely many, and $|W|\le1$ by Lemma~\ref{lem:weilbound}. Since $I_\nu(z)/z^\nu$ is increasing, $I_\nu(z/\gamma)\le\gamma^{-\nu}I_\nu(z)$ for $\gamma\ge1$. So the $\gamma$-th term of \eqref{eq:radmain} is $O(\gamma\cdot\gamma^{-1}\cdot\gamma^{-15/2})$, that of \eqref{eq:radK2} is $O(\gamma^{1-1/2-8})$. In \eqref{eq:radK3}, for $2m\gamma\nmid g$ one has $|\sinh(a-ib)|^2=\sinh^2a+\sin^2b\ge\sin^2b$, so $|f_{\gamma,g}|\le\pi^2\csc^2(\pi g/2m\gamma)$, and for $g=0$ the function $\pi^2/\sinh^2y-\pi^2/y^2$ takes values in $[-\pi^2/3,0)$. With the identity $\sum_{g=1}^{N-1}\csc^2(\pi g/N)=(N^2-1)/3$ for $N=2m\gamma$,
\begin{equation}\label{eq:fsum}
\sum_{g=-m\gamma}^{m\gamma-1}\sup_u|f_{\gamma,g}(u)|\le\frac{\pi^2(4m^2\gamma^2-1)}3+\frac{\pi^2}3=\frac{4\pi^2m^2\gamma^2}3,
\end{equation}
so the $\gamma$-th term of \eqref{eq:radK3} is $O(\gamma\cdot\gamma^{-2}\cdot\gamma^2\cdot\gamma^{-17/2})$. Only absolute bounds are used below, except for the terms with $\gamma=1$ in the frame $S$, which are computed explicitly in Section~\ref{sec:model}.
\end{proof}

\begin{remark}\label{rem:numerics}
The proof of Theorem~\ref{thm:rad} uses no computation; its ingredients were also checked independently. The wall sum \eqref{eq:Pinf} agrees with the coefficients of $\varphi^F_m$ computed exactly from the product formula on all $1498$ polar terms with $m\le12$, and \eqref{eq:P0} agrees with the coefficients of the finite part in the frame $S$, computed exactly from the second product of Section~\ref{sec:polar0}, on all $954$ polar terms with $m\le12$. The right side of \eqref{eq:hypR}, truncated at $\gamma\le8$ in $\mathcal M$ and at $\gamma\le4$ in $\mathcal K_2$ and $\mathcal K_3$, agrees with the exact coefficients for $1\le Q^2\le4$, $n\in\{4,8,16,30\}$ and $0\le\ell\le m$: the relative difference is below $6\cdot10^{-9}$ at $n=4$, $5\cdot10^{-12}$ at $n=8$ and $5\cdot10^{-21}$ at $n=30$, and it decreases with $n$ as the omitted terms with larger $\gamma$ do. The programs are listed in Appendix~\ref{app:interval}.
\end{remark}

\begin{remark}
What the proof of Theorem~\ref{thm:all} uses from (P0) is the structure of Lemma~\ref{lem:walls}: the leading family of walls and the bound (vi) for the others. From (P$\infty$) it uses only that every wall in it has $r$ or $s$ even, hence $rs\ge2$, so that Lemma~\ref{lem:walls}(vi) applies.
\end{remark}
